\documentclass{article}
\usepackage{amsmath,amssymb,amsthm}
\usepackage{bm}
\usepackage{algorithm}
\usepackage[noend]{algpseudocode}
\usepackage{hyperref}
\newtheorem{theorem}{Theorem}
\newtheorem{lemma}{Lemma}
\newtheorem{assumption}{Assumption}
\newcommand{\diam}{\operatorname{diam}}
\newcommand{\gap}{\operatorname{gap}}
\newcommand{\argmin}{\operatorname*{arg\,min}}
\newcommand{\ip}[2]{\langle #1,#2\rangle}
\newcommand{\R}{\mathbb{R}}

\begin{document}

\section*{Algorithm Name}
\textbf{Frank–Wolfe (Conditional Gradient) Method for Convex Smooth Minimization}

\section*{Mathematical Formulation}
Let $f:\R^{d}\rightarrow\R$ be a convex and $L$‑smooth function and let 
\[
\mathcal{C}\subset\R^{d}
\]
be a non‑empty compact convex set (the feasible domain).  
Denote the optimal solution by 
\[
x^{\star}\in\arg\min_{x\in\mathcal{C}}f(x).
\]

At iteration $k\ge 0$ the algorithm performs

\begin{align}
s_{k} &:= \argmin_{s\in\mathcal{C}} \;\ip{\nabla f(x_{k})}{s}, \tag{1}\\[4pt]
\gamma_{k} &\in (0,1] \qquad\text{(step‑size, e.g. $\gamma_{k}=2/(k+2)$ or line‑search)},\\[4pt]
x_{k+1} &:= (1-\gamma_{k})\,x_{k} + \gamma_{k}\,s_{k}. \tag{2}
\end{align}

The linear minimization oracle (LMO) in (1) returns the vertex of $\mathcal{C}$
that minimizes the linearization of $f$ at $x_{k}$.

\section*{Pseudocode}
\begin{algorithm}[H]
\caption{Frank–Wolfe}
\begin{algorithmic}[1]
\State \textbf{Input:} convex $L$‑smooth $f$, compact convex $\mathcal{C}$, 
initial point $x_{0}\in\mathcal{C}$, step‑size rule $\{\gamma_{k}\}$, 
max. iterations $K$.
\For{$k=0,\dots,K-1$}
    \State Compute gradient $g_{k}\gets\nabla f(x_{k})$.
    \State \textbf{LMO:} $s_{k}\gets\argmin_{s\in\mathcal{C}}\ip{g_{k}}{s}$.
    \State Choose step size $\gamma_{k}\in(0,1]$ 
          (e.g. $\gamma_{k}=2/(k+2)$ or exact line search).
    \State Update $x_{k+1}\gets (1-\gamma_{k})x_{k}+\gamma_{k}s_{k}$.
\EndFor
\State \textbf{Output:} $x_{K}$ (or the best iterate).
\end{algorithmic}
\end{algorithm}

\section*{Dry Run Example}
Minimize the univariate quadratic 
\[
f(w)=(w-3)^{2}
\]
over the interval $\mathcal{C}=[0,5]$.  
The gradient is $\nabla f(w)=2(w-3)$.  
We start from $w_{0}=0$ and use a \emph{fixed} step size
$\gamma=0.1$ (i.e. $\gamma_{k}\equiv0.1$).

\begin{center}
\begin{tabular}{c|c|c|c|c}
\textbf{Iter.} $k$ & $w_{k}$ & $\nabla f(w_{k})$ & $s_{k}$ (LMO) & $w_{k+1}$ \\ \hline
0 & $0$ & $2(0-3)=-6$ & $\argmin_{s\in[0,5]}(-6)s = 5$ & $0.9\cdot0+0.1\cdot5=0.5$\\
1 & $0.5$ & $2(0.5-3)=-5$ & $\argmin_{s\in[0,5]}(-5)s = 5$ & $0.9\cdot0.5+0.1\cdot5=0.95$\\
2 & $0.95$ & $2(0.95-3)=-4.1$ & $\argmin_{s\in[0,5]}(-4.1)s = 5$ & $0.9\cdot0.95+0.1\cdot5=1.355$\\
\end{tabular}
\end{center}

Objective values:

\[
f(0)=9,\qquad f(0.5)=6.25,\qquad f(0.95)=4.2025,\qquad f(1.355)\approx3.417.
\]

The iterates move toward the minimizer $w^{\star}=3$ while staying inside $\mathcal{C}$.

\newpage
\section*{Assumptions}
\begin{assumption}[Convexity and Smoothness]\label{ass:convex}
$f:\R^{d}\to\R$ is convex and $L$‑smooth, i.e. for all $x,y\in\R^{d}$
\[
f(y)\le f(x)+\ip{\nabla f(x)}{y-x}+\frac{L}{2}\|y-x\|^{2}. \tag{3}
\]
\end{assumption}
\begin{assumption}[Compact Feasible Set]\label{ass:compact}
$\mathcal{C}\subset\R^{d}$ is convex, compact, and has finite diameter
\[
D:=\diam(\mathcal{C})=\max_{u,v\in\mathcal{C}}\|u-v\|<\infty. \tag{4}
\]
\end{assumption}
\begin{assumption}[Linear Minimization Oracle]\label{ass:lmo}
For every $x\in\mathcal{C}$ the oracle (1) returns an exact solution
$s\in\mathcal{C}$.
\end{assumption}
All subsequent results are derived under \ref{ass:convex}–\ref{ass:lmo}.

\section*{Main Convergence Theorem}
\begin{theorem}[Sublinear $O(1/k)$ rate]\label{thm:fw}
Let $\{x_{k}\}$ be generated by the Frank–Wolfe method with step size
\[
\gamma_{k}=\frac{2}{k+2},\qquad k\ge0. \tag{5}
\]
Define the \emph{Frank–Wolfe gap}
\[
\gap(x_{k})\;:=\;\max_{s\in\mathcal{C}}\ip{\nabla f(x_{k})}{x_{k}-s}
            \;=\;\ip{\nabla f(x_{k})}{x_{k}-s_{k}}. \tag{6}
\]
Then for every $k\ge0$
\[
f(x_{k})-f(x^{\star})\;\le\;\frac{2L D^{2}}{k+2}, \qquad
\gap(x_{k})\;\le\;\frac{4L D^{2}}{k+2}. \tag{7}
\]
\end{theorem}

\section*{Theoretical Properties}
\begin{itemize}
\item \textbf{Stability.} The iterates remain in $\mathcal{C}$ because each update is a convex combination of two points in $\mathcal{C}$.
\item \textbf{Step‑size sensitivity.} The closed‑form schedule (5) yields the optimal $O(1/k)$ bound without line‑search. Any $\gamma_{k}\in[0,1]$ satisfying $\sum_{k}\gamma_{k}=+\infty$ and $\sum_{k}\gamma_{k}^{2}<\infty$ also guarantees convergence, but the constant degrades.
\item \textbf{Comparison.} Compared with projected gradient descent (PGD) which needs a Euclidean projection, Frank–Wolfe replaces the projection by a cheap linear minimization oracle; however the $O(1/k)$ rate is slower than the $O(1/k^{2})$ rate of accelerated PGD for smooth convex problems.
\end{itemize}

\section*{Discussion}
The theorem shows that the Frank–Wolfe gap, a natural certificate of optimality, decays at the same $O(1/k)$ rate as the primal suboptimality. The bound depends only on the smoothness constant $L$ and the geometry of $\mathcal{C}$ through its diameter $D$, making it dimension‑free.

\newpage
\section*{Preliminaries (Definitions and Lemmas)}
\begin{lemma}[Smoothness inequality]\label{lem:smooth}
Under Assumption~\ref{ass:convex}, for any $x,y\in\R^{d}$
\[
f(y)\le f(x)+\ip{\nabla f(x)}{y-x}+\frac{L}{2}\|y-x\|^{2}. \tag{8}
\]
\end{lemma}
\begin{proof}
Directly the statement of Assumption~\ref{ass:convex}. \hfill $\square$
\end{proof}

\begin{lemma}[Curvature constant]\label{lem:curv}
Define the curvature of $f$ over $\mathcal{C}$ as
\[
C_{f}:=\sup_{\substack{x,s\in\mathcal{C}\\\gamma\in[0,1]}}
\frac{2}{\gamma^{2}}\Bigl(f\bigl(x+\gamma(s-x)\bigr)-f(x)-\gamma\ip{\nabla f(x)}{s-x}\Bigr).
\tag{9}
\]
Then $C_{f}\le L D^{2}$. 
\end{lemma}
\begin{proof}
Apply Lemma~\ref{lem:smooth} with $y=x+\gamma(s-x)$:
\[
f(x+\gamma(s-x))\le f(x)+\gamma\ip{\nabla f(x)}{s-x}
                    +\frac{L}{2}\gamma^{2}\|s-x\|^{2}.
\]
Rearranging gives the claimed bound since $\|s-x\|\le D$. \hfill $\square$
\end{proof}

\section*{Main Proof (Step‑by‑Step Derivation)}
We prove Theorem~\ref{thm:fw}. Throughout the proof we number each non‑trivial algebraic manipulation as required.

\subsection*{Step 1 – Descent Lemma for the Update}
From the update rule (2) we have
\[
x_{k+1}=x_{k}+\gamma_{k}(s_{k}-x_{k}). \tag{10}
\]
Applying Lemma~\ref{lem:smooth} with $x=x_{k}$ and $y=x_{k+1}$ yields
\[
f(x_{k+1})\le f(x_{k})+\gamma_{k}\ip{\nabla f(x_{k})}{s_{k}-x_{k}}
            +\frac{L}{2}\gamma_{k}^{2}\|s_{k}-x_{k}\|^{2}. \tag{11}
\]
\textbf{[Step 1]}

\subsection*{Step 2 – Use of the Linear Minimization Oracle}
By definition of $s_{k}$ (1),
\[
\ip{\nabla f(x_{k})}{s_{k}}=\min_{s\in\mathcal{C}}\ip{\nabla f(x_{k})}{s}
\le \ip{\nabla f(x_{k})}{x^{\star}}. \tag{12}
\]
Subtract $\ip{\nabla f(x_{k})}{x_{k}}$ from both sides and rearrange:
\[
\ip{\nabla f(x_{k})}{s_{k}-x_{k}}
\le \ip{\nabla f(x_{k})}{x^{\star}-x_{k}}
= -\gap(x_{k}). \tag{13}
\]
\textbf{[Step 2]}

\subsection*{Step 3 – Bounding the Distance Term}
Because $x_{k},s_{k}\in\mathcal{C}$, the diameter bound (4) gives
\[
\|s_{k}-x_{k}\|\le D. \tag{14}
\]
\textbf{[Step 3]}

\subsection*{Step 4 – Combine (11)–(14)}
Insert (13) and (14) into (11):
\[
f(x_{k+1})\le f(x_{k})-\gamma_{k}\gap(x_{k})
            +\frac{L}{2}\gamma_{k}^{2}D^{2}. \tag{15}
\]
\textbf{[Step 4]}

\subsection*{Step 5 – Relate Gap to Suboptimality}
Convexity of $f$ yields for any $x\in\mathcal{C}$
\[
f(x)-f(x^{\star})\le \ip{\nabla f(x)}{x-x^{\star}}. \tag{16}
\]
Taking $x=x_{k}$ and using the definition (6) of the gap,
\[
\gap(x_{k})=\ip{\nabla f(x_{k})}{x_{k}-s_{k}}
           \ge \ip{\nabla f(x_{k})}{x_{k}-x^{\star}}
           = f(x_{k})-f(x^{\star}). \tag{17}
\]
\textbf{[Step 5]}

\subsection*{Step 6 – Recurrence on Suboptimality}
Define $\Delta_{k}:=f(x_{k})-f(x^{\star})\ge0$.  
Using (15) and the lower bound (17):
\[
\Delta_{k+1}\le \Delta_{k}-\gamma_{k}\Delta_{k}
               +\frac{L}{2}\gamma_{k}^{2}D^{2}
            = (1-\gamma_{k})\Delta_{k}
              +\frac{L}{2}\gamma_{k}^{2}D^{2}. \tag{18}
\]
\textbf{[Step 6]}

\subsection*{Step 7 – Plug the Specific Step Size}
Set $\gamma_{k}=2/(k+2)$ as in (5). Then
\[
1-\gamma_{k}=1-\frac{2}{k+2}=\frac{k}{k+2},\qquad
\gamma_{k}^{2}= \frac{4}{(k+2)^{2}}. \tag{19}
\]
Insert (19) into (18):
\[
\Delta_{k+1}\le \frac{k}{k+2}\Delta_{k}
               +\frac{L D^{2}}{(k+2)^{2}}. \tag{20}
\]
\textbf{[Step 7]}

\subsection*{Step 8 – Induction to Obtain Closed Form}
We claim by induction that for all $k\ge0$
\[
\Delta_{k}\le \frac{2L D^{2}}{k+2}. \tag{21}
\]
\textbf{Base $k=0$:}  
$\Delta_{0}=f(x_{0})-f(x^{\star})\le L D^{2}$ (by $L$‑smoothness and $x_{0},x^{\star}\in\mathcal{C}$), which is certainly $\le 2L D^{2}/2$.

\textbf{Inductive step:} Assume (21) holds for $k$. Using (20):
\[
\Delta_{k+1}\le \frac{k}{k+2}\cdot\frac{2L D^{2}}{k+2}
               +\frac{L D^{2}}{(k+2)^{2}}
            =\frac{2L D^{2}\,k + L D^{2}}{(k+2)^{2}}
            =\frac{L D^{2}(2k+1)}{(k+2)^{2}}.
\]
Observe that
\[
\frac{L D^{2}(2k+1)}{(k+2)^{2}}
   \le \frac{2L D^{2}}{k+3},
\]
because the inequality reduces to $(2k+1)(k+3)\le 2(k+2)^{2}$, which holds for all $k\ge0$. Hence (21) holds for $k+1$.
\textbf{[Step 8]}

\subsection*{Step 9 – Gap Bound}
From (17) and (21) we obtain
\[
\gap(x_{k})\le \Delta_{k}+\frac{L}{2}\gamma_{k}^{2}D^{2}
            \le \frac{2L D^{2}}{k+2}+\frac{L D^{2}}{(k+2)^{2}}
            \le \frac{4L D^{2}}{k+2}. \tag{22}
\]
The last inequality uses $1/(k+2)^{2}\le 1/(k+2)$.
\textbf{[Step 9]}

\subsection*{Conclusion}
Inequalities (21) and (22) are exactly the statements in \eqref{thm:fw}. ∎

\section*{Rate Derivation (With Constants)}
The constant $2L D^{2}$ in (21) comes from the curvature bound $C_{f}\le L D^{2}$ (Lemma~\ref{lem:curv}) and the choice $\gamma_{k}=2/(k+2)$. If a different step‑size rule $\gamma_{k}=2/(k+\tau)$ with $\tau\ge2$ is used, the bound becomes $C_{f}/(\tau-1)$.

\section*{Special Cases}
\begin{itemize}
\item \textbf{Polyhedral $\mathcal{C}$.} When $\mathcal{C}$ is a polytope, the LMO reduces to solving a linear program; the above rates still hold.
\item \textbf{Strongly convex $f$.} If $f$ is $\mu$‑strongly convex, the primal gap decays linearly after a finite number of iterations (see \cite{Jaggi2013}).
\item \textbf{Exact line search.} Choosing $\gamma_{k}=\argmin_{\gamma\in[0,1]}f(x_{k}+\gamma(s_{k}-x_{k}))$ yields the same $O(1/k)$ rate with a possibly smaller constant.
\end{itemize}

\begin{thebibliography}{9}
\bibitem{Jaggi2013}
M.~Jaggi, \emph{Revisiting Frank–Wolfe: Projection‑Free Sparse Convex Optimization}, Proc. ICML, 2013.
\end{thebibliography}

\end{document}